\chapter{El modelo de planeador}\label{ch:planeador}
\objetivos{Derivar el equilibrio de planeo; relacionar polar, ángulo y velocidades; distinguir relaciones exactas de aproximaciones de pequeño ángulo.}

\section{La energía que sostiene el movimiento}
Un planeador en aire en calma consume energía potencial para compensar el trabajo de la resistencia. Puede mantener velocidad constante mientras desciende porque la componente del peso paralela a la trayectoria equilibra $D$. Si además mantiene altura y no recibe energía del aire o de una planta propulsora, no puede sostener un vuelo estacionario con resistencia positiva.

Se adopta $\gamma<0$ en descenso y empuje nulo. Las ecuaciones de punto material son
\begin{equation}
m\dot V=-D-W\sin\gamma,\qquad mV\dot\gamma=L-W\cos\gamma.
\end{equation}
En planeo rectilíneo estacionario,
\begin{equation}
D=-W\sin\gamma,\quad L=W\cos\gamma,\quad
\tan(-\gamma)=D/L=C_D/C_L.\label{eq:planeo-exacto}
\end{equation}
La razón $L/D$ es la eficiencia aerodinámica. En calma coincide con el avance horizontal por unidad de altura perdida, porque $\dot x/(-\dot h)=\cot(-\gamma)=L/D$. Esta igualdad geométrica es exacta para el movimiento definido.

\ilustra{fuerzas-planeo}{Equilibrio de un planeo estacionario: sustentación normal a la velocidad aérea, resistencia opuesta a ella y peso vertical. Los tres vectores de fuerza están dibujados con una escala común y cierran el equilibrio; la velocidad usa una escala independiente.}

\section{Velocidad y descenso exactos}
Definiendo $\delta=-\gamma>0$, para un $C_L$ elegido se calcula $\delta=\arctan(C_D/C_L)$. La velocidad necesaria es
\begin{equation}
V=\sqrt{\frac{2W\cos\delta}{\rho SC_L}},\qquad
s=-\dot h=V\sin\delta,\qquad V_H=V\cos\delta.
\label{eq:planeo-velocidad}
\end{equation}
Combinando $\cos\delta=C_L/\sqrt{C_L^2+C_D^2}$ se obtiene otra forma útil:
\begin{equation}
V=\sqrt{\frac{2W}{\rho S}}(C_L^2+C_D^2)^{-1/4},\quad
s=\sqrt{\frac{2W}{\rho S}}\frac{C_D}{(C_L^2+C_D^2)^{3/4}}.
\end{equation}
Las dos expresiones de velocidad deben coincidir; sirven como comprobación independiente de implementación. La velocidad horizontal no es TAS y la tasa de descenso no es el ángulo de planeo.

\section{Aproximación de planeo tendido}
Para $C_D/C_L\ll1$, $\cos\delta\simeq1$ y $\sin\delta\simeq\tan\delta\simeq C_D/C_L$. Así,
\begin{equation}
V\simeq\sqrt{\frac{2W}{\rho SC_L}},\qquad
s\simeq\sqrt{\frac{2W}{\rho S}}\frac{C_D}{C_L^{3/2}}.
\label{eq:descenso-aprox}
\end{equation}
La aproximación $L\simeq W$ simplifica la optimización, pero no debe mezclarse con una afirmación de exactitud en descensos pronunciados. Con eficiencia 25, $\delta\simeq2.29^\circ$ y el error relativo de sustituir $\cos\delta$ por uno es aproximadamente $0.08\%$. Con eficiencia 3, el error ya supera el 5\%.

\section{Polar de descenso}
Al recorrer los $C_L$ admisibles se obtiene una curva de tasa de descenso frente a velocidad. En la aproximación tendida, la tangente desde el origen identifica la mínima razón $s/V$, relacionada con mejor planeo. El punto más bajo identifica mínimo descenso. Son criterios diferentes: uno maximiza distancia por altura y otro tiempo por altura.

\ilustra{polar-descenso}{Polar de descenso de un planeador docente. Se emplea la aproximación tendida con masa constante y densidad de nivel del mar. El dominio de utilización debe recortarse con la pérdida y otros límites.}

El límite $V\to0$ de la expresión parabólica no describe un vuelo posible, pues exige $C_L\to\infty$. En el extremo rápido, una polar de baja velocidad tampoco describe necesariamente el aumento de resistencia por compresibilidad o los límites estructurales. Los extremos matemáticos no eliminan la necesidad de una envolvente.

\section{Influencia de la masa y de la densidad}
A igual $C_L$ y polar, $\delta$ no depende del peso. Tanto $V$ como $s$ varían como $\sqrt{W/\rho}$. Por eso un planeador lastrado puede conservar aproximadamente el mismo ángulo de planeo y recorrer la misma distancia en calma desde una altura dada, pero lo hace más rápido y permanece menos tiempo en el aire. La sustentación máxima, el número de Reynolds y el estado superficial pueden cambiar la polar real: la independencia del alcance respecto a la masa pertenece al modelo, no a cualquier avión en cualquier condición.

A mayor altura ISA, la TAS de un $C_L$ dado aumenta. La EAS correspondiente es aproximadamente constante para masa constante en planeo tendido. No se debe introducir la densidad de la altitud inicial en toda una integral de descenso de varios kilómetros sin evaluar el error.

\caso{Ejemplo original: un punto de planeo}{Sean $m=420\SIu{kg}$, $S=15.8\SIu{m^2}$, $C_{D0}=0.013$, $k=0.031$, $C_L=0.65$ y $\rho=1.225\SIu{kg\,m^{-3}}$. Entonces $C_D=0.02610$, $L/D=24.91$ y $\gamma\simeq-2.30^\circ$. La velocidad se calcula con $\cos\gamma$ en el numerador; la tasa de descenso se obtiene como $-V\sin\gamma$, no multiplicando por un ángulo en grados. Desde $1000\SIu m$ de altura sobre el punto final, el alcance geométrico en calma es aproximadamente $24.91\SIu{km}$, si el mismo $C_L$ y polar se mantienen.}

\section{Balance energético como segunda vía}
Multiplicando la ecuación tangencial por $V$ y utilizando $\dot h=V\sin\gamma$,
\begin{equation}
\frac{\dd}{\dd t}\left(\tfrac12mV^2+mgh\right)=-DV.
\end{equation}
En planeo estacionario, $mgs=DV$ y $s=DV/W$. Esta relación es exacta cuando la velocidad es constante. Reemplazar además $L$ por $W$ al evaluar $D$ introduce la aproximación de ángulo pequeño. Una recuperación de velocidad después de una perturbación puede intercambiar energía cinética y potencial; no se predice con una simple tabla de descenso estacionario.

\section{Taller resuelto: Medir el error de planeo tendido}

\noindent\textbf{Nivel:} Intermedio.

\apartado{1. Enunciado}{Un planeador de 455 kg, S=16.2 m² y polar 0.014+0.032CL² vuela a CL=0.78 en rho=1.05. Compare TAS y descenso exactos con la aproximación L=W y seno de delta igual a tangente.}

\apartado{2. Datos}{Equilibrio estacionario sin empuje; gamma negativo.}

\apartado{3. Incógnitas}{Las magnitudes y comprobaciones solicitadas en el caso.}

\apartado{4. Sistema de referencia}{Trayectoria aérea en calma y fuerzas viento.}

\apartado{5. Hipótesis físicas}{La misma polar se utiliza en ambas comparaciones; sólo cambian las aproximaciones trigonométricas.}

\apartado{6. Modelo matemático}{Medir el error de planeo tendido}

\apartado{7. Ecuaciones simbólicas}{Se parte de los balances y definiciones siguientes.}

\begin{equation}\delta=\arctan(C_D/C_L)\end{equation}

\begin{equation}V_{exacto}=\sqrt{2W\cos\delta/(\rho SC_L)}\end{equation}

\begin{equation}s_{exacto}=V_{exacto}\sin\delta\end{equation}

\begin{equation}V_{aprox}=\sqrt{2W/(\rho SC_L)},\quad s_{aprox}=V_{aprox}C_D/C_L\end{equation}

\apartado{8. Desarrollo algebraico}{Las operaciones se realizan en este orden, conservando las condiciones del modelo.}

\begin{enumerate}

\item Se evalúa un único CD y ángulo para ambas vías.

\item Se conserva coseno en el equilibrio exacto.

\item Se calcula el error relativo de descenso y se interpreta su tamaño respecto a datos de polar.

\end{enumerate}

\apartado{9. Sustitución numérica}{Se conservan los resultados intermedios sin redondear y se muestran todas las operaciones escalares disponibles.}

\textbf{Desarrollo de las magnitudes intermedias.} De $D=W\sin\delta$ y $L=W\cos\delta$, dividir da $\tan\delta=C_D/C_L$. La sustentación exacta usa el coseno; el descenso usa el seno. La aproximación tendida sustituye ambos factores por sus límites de ángulo pequeño.\begin{enumerate}

\item Calcular la polar al CL impuesto.

\[\fitmath{C_D=0.014+0.032\cdot {0.78}^{2}\approx 0.0334688\,}\]

\item Obtener ángulo de descenso positivo.

\[\fitmath{\delta=\arctan\left(\frac{C_D}{0.78}\right)}\]

\[\fitmath{\delta=\arctan\left(\frac{(0.0334688)}{0.78}\right)\approx 0.04288241\,\mathrm{rad}}\]

\item Conservar coseno en el equilibrio normal.

\[\fitmath{\mathrm{Vexact}=\sqrt{\frac{2\cdot 455\cdot 9.80665\cdot \cos\left(\delta\right)}{1.05\cdot 16.2\cdot 0.78}}}\]

\[\fitmath{\mathrm{Vexact}=\sqrt{\frac{2\cdot 455\cdot 9.80665\cdot \cos\left((0.04288241309)\right)}{1.05\cdot 16.2\cdot 0.78}}\approx 25.9228\,\mathrm{m/s}}\]

\item Proyectar la velocidad sobre la vertical.

\[\fitmath{\mathrm{sexact}=\mathrm{Vexact}\cdot \sin\left(\delta\right)}\]

\[\fitmath{\mathrm{sexact}=(25.92279697)\cdot \sin\left((0.04288241309)\right)\approx 1.111291\,\mathrm{m/s}}\]

\item Aplicar solo en esta ruta el equilibrio aproximado L=W.

\[\fitmath{\mathrm{Vapprox}=\sqrt{\frac{2\cdot 455\cdot 9.80665}{1.05\cdot 16.2\cdot 0.78}}\approx 25.93472\,\mathrm{m/s}}\]

\item Sustituir seno por tangente en la aproximación.

\[\fitmath{\mathrm{sapprox}=\frac{\mathrm{Vapprox}\cdot C_D}{0.78}}\]

\[\fitmath{\mathrm{sapprox}=\frac{(25.93472073)\cdot (0.0334688)}{0.78}\approx 1.112826\,\mathrm{m/s}}\]

\item Cuantificar el error en descenso.

\[\fitmath{\mathrm{relative}=\frac{\mathrm{sapprox}}{\mathrm{sexact}}-1}\]

\[\fitmath{\mathrm{relative}=\frac{(1.112825617)}{(1.111291422)}-1\approx 0.001380551\,}\]

\end{enumerate}





\apartado{10. Resultados con unidades}{Los resultados pertenecen al modelo y al intervalo de datos declarados.}

\begin{center}\begin{tabular}{lr}\toprule Magnitud&Resultado\\\midrule

TAS exacta & $25.9228\,\mathrm{m\,s^{-1}}$\\

Descenso exacto & $1.111291\,\mathrm{m\,s^{-1}}$\\

Descenso aproximado & $1.112826\,\mathrm{m\,s^{-1}}$\\

Error relativo de descenso & $0.001380551\,$\\

\bottomrule\end{tabular}\end{center}

\apartado{11. Interpretación}{La aproximación es buena para una eficiencia elevada, pero la igualdad exacta y la aproximada no son la misma relación.}

\apartado{12. Comprobación independiente y alcance}{Comprobar que el cociente de velocidades es sqrt(cos(delta)) y evaluar energía DV/W en el caso exacto.}

\apartado{13. Limitaciones}{Un error pequeño de trigonometría no compensa incertidumbre grande de la polar o falta de trim.}

\par\smallskip\noindent{\footnotesize\textit{Datos y solución originales. Las sustituciones registradas se reproducen con SymPy y Studium; eso no establece validación experimental ni revisión experta.}}


\repaso{¿Mejora la eficiencia aerodinámica al añadir lastre si la polar permanece igual?}{No en este modelo. Cambian las velocidades características y la tasa de descenso. El beneficio de penetración frente a viento depende de la velocidad sobre el suelo y de las condiciones de la misión.}
\resumen{El planeo equilibra resistencia con peso tangencial y sustentación con peso normal. El alcance por altura depende de $L/D$, mientras que el tiempo depende también de la velocidad. La aproximación tendida se usa con su error y su dominio declarados.}
\fuentes{How, MIT 16.333, lección 1, secciones de planeo y balance energético. Familias de ejercicios identificadas en los materiales facilitados; datos y desarrollos de esta edición son originales.}
